% ============================================================
% 02_resumen_paper.tex
% Responsable: Integrante X
% ============================================================

\section{Resumen del Paper}

\subsection{Problema abordado}

% TODO: Describir el problema exacto que resuelve C-Lock.
% - Overhead energético de busy-waiting en spin-locks
% - El procesador no puede apagarse mientras espera el lock

\subsection{Propuesta: C-Lock}

C-Lock es un mecanismo hardware de sincronización que utiliza \textit{clock-gating} para
detener el reloj de los cores que esperan un lock, reduciendo el consumo energético sin
degradar el rendimiento.

La idea central es que, en lugar de que un core ejecute instrucciones en bucle esperando
un lock, el hardware detecta el conflicto y desactiva el reloj del core hasta que el
recurso esté disponible.

\subsection{Componentes principales}

\begin{itemize}
    \item \textbf{C-LockManager:} Componente central. Contiene el árbitro, el contador
          global y los \textit{pools}.
    \item \textbf{Árbitro (Arbiter):} Gestiona el orden de acceso a los locks.
    \item \textbf{Contador global (Global Counter):} Registro compartido para tracking
          de locks activos.
    \item \textbf{Pool\_i:} Un pool por core, contiene los \textit{Conflict Checkers}.
    \item \textbf{Conflict Checker:} Detecta si existe conflicto para un lock dado y
          activa el clock-gating si es necesario.
\end{itemize}

\subsection{Interfaz software-hardware}

El programador usa tres macros para interactuar con C-Lock:

\begin{lstlisting}[language=C, caption={Macros de C-Lock}]
ADD_ITEM(lock_id);        // Registrar lock al inicio del programa
BEGIN_C_LOCK(lock_id);    // Adquirir lock (puede gatillar clock-gating)
END_C_LOCK(lock_id);      // Liberar lock
\end{lstlisting}

\subsection{Evaluación y resultados}

% TODO: Completar con los resultados del paper.
% - Benchmarks usados: STAMP, MiBench, TSM
% - Reducción de consumo energético reportada
% - Impacto en rendimiento (speedup / overhead)

Los autores evaluaron C-Lock usando los benchmarks \textit{STAMP}, \textit{MiBench}
y \textit{TSM}, reportando una reducción significativa en el consumo energético
comparado con spin-locks clásicos, con un impacto mínimo en el rendimiento.
